Un ensayo simulado representa un depósito bien mezclado de \(100\,\mathrm L\). Entra agua limpia a caudal constante desconocido \(q>0\), en \(\mathrm{L/min}\), y sale mezcla al mismo caudal. Además, un proceso interno consume por minuto una fracción \(k\ge0\) de la masa de contaminante presente; \(k\) se mide en \(\mathrm{min}^{-1}\). La concentración registrada, tratada como exacta en este ejercicio, es
\[
c(t)=2e^{-0.12t}\ \mathrm{g/L},\qquad t\ge0.
\]
Formule el balance y determine todo lo que esta curva permite inferir sobre \(q\) y \(k\). Se puede repetir el ensayo sin cambiar el régimen hidráulico, añadiendo una pequeña cantidad de un trazador que no reacciona, no se adsorbe y no altera el proceso. ¿Aporta más información volver a medir el volumen constante o medir la disminución del trazador? En la repetición, la concentración del trazador se reduce a la mitad en diez minutos. Determine entonces \(q\) y \(k\), y compruebe su compatibilidad con las hipótesis.
